% ============================================================
%  Flipped UPQC for MV/LV Distribution Networks
%  PCIM Europe Conference Paper
%  Author: Davide Bagnara
% ============================================================

\documentclass[conference]{IEEEtran}

% --- Encoding & Language ---
\usepackage[T1]{fontenc}
\usepackage[utf8]{inputenc}
\usepackage[english]{babel}

% --- Math ---
\usepackage{amsmath}
\usepackage{amssymb}
\usepackage{bm}          % bold math symbols

% --- Graphics & Figures ---
\usepackage{graphicx}
\usepackage{float}
\usepackage{subfig}

% --- Tables ---
\usepackage{booktabs}
\usepackage{multirow}

% --- Circuit / Block Diagrams ---
\usepackage{tikz}
\usepackage{circuitikz}
\usetikzlibrary{arrows.meta, positioning, shapes.geometric, calc}

% --- Hyperlinks (load last) ---
\usepackage[hidelinks]{hyperref}

% ============================================================
%  TITLE & AUTHORS
% ============================================================

\title{%
  A Flipped Unified Power Quality Conditioner\\
  for Active Distribution Networks:\\
  Topology, Modelling, and Control Strategy%
}

\author{%
  \IEEEauthorblockN{Davide Bagnara}
  \IEEEauthorblockA{%
    Independent Researcher\\
    Sterzing (BZ), Italy\\
    \href{mailto:davide.bagnara@example.com}{davide.bagnara@example.com}%
  }
}

\begin{document}

\maketitle

% ============================================================
%  ABSTRACT
% ============================================================

\begin{abstract}
The increasing penetration of distributed renewable energy resources
(DERs) into medium-voltage (MV) and low-voltage (LV) distribution
networks introduces new challenges for both power quality and grid
stability. Conventional mitigation devices---dynamic voltage restorers
(DVRs), static synchronous compensators (STATCOMs), and
uninterruptible power supplies (UPS)---address these issues in
isolation and adopt a purely reactive role.
This paper proposes a \emph{Flipped Unified Power Quality Conditioner}
(F-UPQC), an architecture in which the shunt converter faces the grid
to provide ancillary services to the distribution system operator
(DSO), while the series converter faces the load to guarantee voltage
stability for sensitive industrial loads.
A shared DC link with optional energy storage completes the topology,
enabling participation in local flexibility markets without additional
hardware.
The mathematical model of the complete system is derived, and a
decoupled control strategy is presented. Theoretical analysis
demonstrates the ability of the proposed topology to simultaneously
compensate voltage sags and swells at the load terminals while
injecting reactive power and mitigating current harmonics on the grid
side. The proposed architecture is aligned with the evolving regulatory
framework of the European Clean Energy Package and the growing demand
for grid-edge devices capable of active distribution network (ADN)
support.
\end{abstract}

\begin{IEEEkeywords}
power quality, unified power quality conditioner, UPQC, active
distribution network, ancillary services, dynamic voltage restorer,
STATCOM, energy storage, voltage sag compensation, reactive power
control
\end{IEEEkeywords}

% ============================================================
%  I. INTRODUCTION
% ============================================================

\section{Introduction}
\label{sec:introduction}

The decarbonisation of the European energy system is driving a rapid
and large-scale deployment of distributed energy resources (DERs)
at the medium-voltage (MV) and low-voltage (LV) distribution level.
Photovoltaic plants, wind micro-turbines, battery storage systems, and
electric vehicle chargers are transforming distribution grids from
passive, unidirectional infrastructures into dynamic, bidirectional
active distribution networks (ADNs)~\cite{ref:EDSO_transition}.
The European Commission has estimated investment requirements of
375--425\,billion EUR in distribution infrastructure by 2030, an
increase of 50--70\,\% compared to the previous decade~\cite{ref:EC_grids}.

This structural transformation carries two simultaneous and
partially conflicting challenges. On one hand, the proliferation of
power-electronics-interfaced generation and storage introduces
harmonics, voltage unbalance, and rapid voltage fluctuations that
degrade power quality for sensitive industrial loads~\cite{ref:IEEE_std_519}.
On the other hand, the reduction of synchronous generation reduces
system inertia and reactive power availability, placing new burdens
on distribution system operators (DSOs) to procure ancillary services
at the distribution level~\cite{ref:ancillary_EU}.
Industry studies estimate that power quality disturbances cost European
and North American industries in excess of 150\,billion USD annually in
lost productivity and equipment damage~\cite{ref:PQ_cost}.

Existing mitigation technologies address these challenges separately.
Dynamic Voltage Restorers (DVRs)~\cite{ref:DVR_original} inject a
compensating series voltage to protect sensitive loads from sags and
swells, but provide no service to the grid.
Static Synchronous Compensators (STATCOMs)~\cite{ref:STATCOM_ref} and
Active Power Filters (APFs)~\cite{ref:APF_ref} operate on the shunt
path and can provide reactive power support or harmonic mitigation,
but offer no load-side protection.
The Unified Power Quality Conditioner (UPQC), first proposed
in~\cite{ref:UPQC_original}, combines both functions by coupling a
shunt and a series converter through a common DC link.
However, in the classical UPQC configuration, the shunt converter is
placed on the \emph{load side} to compensate load current harmonics,
while the series converter is placed on the \emph{grid side} to
compensate supply voltage disturbances.
This arrangement prioritises the perspective of the end user and
provides limited benefit to the DSO.

The present paper proposes a \emph{topological inversion} of the
standard UPQC architecture---hereafter referred to as the
\emph{Flipped UPQC} (F-UPQC). In the proposed configuration:
\begin{itemize}
  \item the \textbf{shunt converter} is connected on the \emph{grid
        side}, enabling provision of reactive power support, harmonic
        current injection, and fast voltage regulation as ancillary
        services to the DSO;
  \item the \textbf{series converter} is connected on the \emph{load
        side}, ensuring full voltage stability for connected sensitive
        loads regardless of grid disturbances;
  \item an \textbf{energy storage} element on the shared DC bus
        decouples active power exchange between the two converters
        and enables participation in distribution-level flexibility
        markets.
\end{itemize}

The main contributions of this paper are:
\begin{enumerate}
  \item A novel F-UPQC topology specifically designed for MV/LV
        distribution networks operating in the context of ADNs;
  \item A complete small-signal mathematical model of the F-UPQC in
        the synchronous reference frame ($dq$ frame);
  \item A decoupled control architecture separating the series
        voltage-stabilisation loop from the shunt ancillary-service
        loop;
  \item A steady-state power flow analysis demonstrating the
        feasibility of simultaneous load protection and DSO service
        provision;
  \item A discussion of the regulatory context (EU Clean Energy
        Package, EN~50160, IEC~61000) and the positioning of the
        F-UPQC in emerging DSO flexibility markets.
\end{enumerate}

The remainder of the paper is organised as follows.
Section~\ref{sec:topology} describes the F-UPQC topology and
compares it with the classical UPQC.
Section~\ref{sec:model} derives the mathematical model.
Section~\ref{sec:control} presents the control strategy.
Section~\ref{sec:steadystate} analyses the steady-state power flow
and component ratings.
Section~\ref{sec:regulatory} discusses regulatory alignment.
Section~\ref{sec:conclusion} draws conclusions and outlines future work.

% ============================================================
%  II. SYSTEM TOPOLOGY
% ============================================================

\section{F-UPQC Topology}
\label{sec:topology}

\subsection{Classical UPQC vs.\ Flipped UPQC}

Fig.~\ref{fig:comparison} illustrates the structural difference between
the classical UPQC and the proposed F-UPQC.
In the classical topology, the shunt converter compensates non-linear
load currents on the load bus, while the series converter mitigates
supply-side disturbances.
In the F-UPQC, the two converters exchange roles with respect to their
point of connection: the shunt VSC interfaces with the grid at the
point of common coupling (PCC), and the series VSC injects a
compensating voltage on the load feeder through a series injection
transformer.

% --- PLACEHOLDER: insert comparison figure ---
\begin{figure}[htbp]
  \centering
  % \includegraphics[width=\columnwidth]{fig/topology_comparison.pdf}
  \fbox{\parbox{0.9\columnwidth}{\centering\vspace{1.5cm}%
    [PLACEHOLDER: Classical UPQC vs.\ F-UPQC block diagram]\vspace{1.5cm}}}
  \caption{Structural comparison between (a) classical UPQC and
           (b) proposed Flipped UPQC (F-UPQC).}
  \label{fig:comparison}
\end{figure}

\subsection{Component Description}

The F-UPQC consists of the following main components:

\subsubsection{Grid-side coupling transformer}
A medium-voltage transformer with Dyn11 vector group provides
galvanic isolation and a 30\textdegree{} phase shift that reduces
low-order current harmonics reflected on the MV bus.

\subsubsection{Shunt VSC (grid-side)}
A two-level or three-level voltage-source converter connected in
shunt at the PCC via a coupling reactor $L_\text{sh}$.
Operating as a grid-following controlled current source, it injects
reactive current $i_q$ on demand from the DSO and a harmonic
cancellation current $i_h$ computed from the grid current measurement.

\subsubsection{Series VSC and injection transformer}
A voltage-source converter whose AC terminals are connected to
the primary of a series injection transformer~$T_\text{s}$.
The secondary winding is inserted in series on the load feeder.
The converter synthesises a compensating voltage
$\bm{v}_\text{inj}$ such that the load terminal voltage
$\bm{v}_\text{load}$ tracks a clean sinusoidal reference
$\bm{v}_\text{ref}$ independently of grid disturbances.

\subsubsection{Common DC bus}
A capacitor bank $C_\text{dc}$ maintains the DC link voltage
$V_\text{dc}$ at a regulated level. The shunt VSC is responsible for
DC bus voltage regulation in normal operation.

\subsubsection{Energy storage interface}
An optional bidirectional DC/DC converter connects an energy storage
system (battery, supercapacitor, or hybrid) to the DC bus.
During sustained voltage sags, the storage supplies the active power
demanded by the series VSC, extending compensation capability
beyond the steady-state limit dictated by grid power flow alone.

% ============================================================
%  III. MATHEMATICAL MODEL
% ============================================================

\section{Mathematical Model}
\label{sec:model}

\subsection{Reference Frame and Notation}

All quantities are expressed in the synchronous $dq$ reference frame
rotating at the fundamental angular frequency $\omega_0 = 2\pi f_0$,
with $f_0 = 50$\,Hz.
The Park transformation adopted is power-invariant:
\begin{equation}
  \bm{T}(\theta) =
  \sqrt{\frac{2}{3}}
  \begin{bmatrix}
    \cos\theta & \cos\!\left(\theta - \tfrac{2\pi}{3}\right)
              & \cos\!\left(\theta + \tfrac{2\pi}{3}\right) \\[4pt]
    -\sin\theta & -\sin\!\left(\theta - \tfrac{2\pi}{3}\right)
               & -\sin\!\left(\theta + \tfrac{2\pi}{3}\right)
  \end{bmatrix}
  \label{eq:park}
\end{equation}
where $\theta = \omega_0 t + \theta_0$ is estimated by a
synchronous-reference-frame phase-locked loop (SRF-PLL).

\subsection{Shunt VSC Model}

Let $\bm{v}_\text{pcc} = [v_{d,\text{pcc}},\, v_{q,\text{pcc}}]^\top$
be the PCC voltage and $\bm{i}_\text{sh} = [i_{d,\text{sh}},\, i_{q,\text{sh}}]^\top$
the shunt converter output current, both in the $dq$ frame.
The coupling reactor dynamics are:
\begin{equation}
  L_\text{sh}\frac{d}{dt}
  \begin{bmatrix} i_{d,\text{sh}} \\ i_{q,\text{sh}} \end{bmatrix}
  =
  \begin{bmatrix} v_{d,\text{sh}} \\ v_{q,\text{sh}} \end{bmatrix}
  -
  \begin{bmatrix} v_{d,\text{pcc}} \\ v_{q,\text{pcc}} \end{bmatrix}
  -
  R_\text{sh}
  \begin{bmatrix} i_{d,\text{sh}} \\ i_{q,\text{sh}} \end{bmatrix}
  +
  \omega_0 L_\text{sh}
  \begin{bmatrix} \phantom{-}i_{q,\text{sh}} \\ -i_{d,\text{sh}} \end{bmatrix}
  \label{eq:shunt_dq}
\end{equation}
where $\bm{v}_\text{sh}$ is the converter output voltage and
$R_\text{sh}$ represents switching and resistive losses.

\subsection{Series VSC Model}

Let $n_s$ be the turns ratio of the series injection transformer
and $\bm{v}_\text{inj}$ the injected voltage referred to the secondary.
The load terminal voltage is:
\begin{equation}
  \bm{v}_\text{load} = \bm{v}_\text{pcc} + \bm{v}_\text{inj}
                      - Z_\text{line}\,\bm{i}_\text{load}
  \label{eq:series_voltage}
\end{equation}
The series converter dynamics in the $dq$ frame follow:
\begin{equation}
  L_\text{se}\frac{d}{dt}
  \begin{bmatrix} i_{d,\text{se}} \\ i_{q,\text{se}} \end{bmatrix}
  =
  \begin{bmatrix} v_{d,\text{se}} \\ v_{q,\text{se}} \end{bmatrix}
  - n_s
  \begin{bmatrix} v_{d,\text{inj}} \\ v_{q,\text{inj}} \end{bmatrix}
  - R_\text{se}
  \begin{bmatrix} i_{d,\text{se}} \\ i_{q,\text{se}} \end{bmatrix}
  + \omega_0 L_\text{se}
  \begin{bmatrix} \phantom{-}i_{q,\text{se}} \\ -i_{d,\text{se}} \end{bmatrix}
  \label{eq:series_dq}
\end{equation}

\subsection{DC Link Model}

The DC bus energy balance is:
\begin{equation}
  C_\text{dc}\frac{dV_\text{dc}}{dt}
  = i_\text{sh,dc} - i_\text{se,dc} + i_\text{sto}
  \label{eq:dclink}
\end{equation}
where $i_\text{sh,dc}$ and $i_\text{se,dc}$ are the DC currents
drawn by the shunt and series VSCs respectively, and $i_\text{sto}$
is the storage current (positive when discharging).
In the power-invariant $dq$ frame:
\begin{align}
  i_\text{sh,dc} &= \frac{v_{d,\text{sh}}\,i_{d,\text{sh}}
                        + v_{q,\text{sh}}\,i_{q,\text{sh}}}{V_\text{dc}}
  \label{eq:psh_dc} \\
  i_\text{se,dc} &= \frac{v_{d,\text{se}}\,i_{d,\text{se}}
                        + v_{q,\text{se}}\,i_{q,\text{se}}}{V_\text{dc}}
  \label{eq:pse_dc}
\end{align}

% ============================================================
%  IV. CONTROL STRATEGY
% ============================================================

\section{Control Strategy}
\label{sec:control}

\subsection{Overview}

The F-UPQC control is decomposed into four decoupled loops, as shown
in Fig.~\ref{fig:control_overview}:
\begin{enumerate}
  \item \textbf{Shunt current controller} --- inner PI loop regulating
        $i_{d,\text{sh}}$ and $i_{q,\text{sh}}$ with cross-coupling
        feed-forward decoupling;
  \item \textbf{DC bus voltage controller} --- outer PI loop providing
        the $d$-axis current reference $i_{d,\text{sh}}^*$ to the shunt
        inner loop, ensuring DC link energy balance;
  \item \textbf{Ancillary service controller} --- generates the
        $q$-axis reference $i_{q,\text{sh}}^*$ according to the
        reactive power demand $Q^*$ dispatched by the DSO;
  \item \textbf{Series voltage controller} --- computes the required
        injection voltage $\bm{v}_\text{inj}^*$ to track
        $\bm{v}_\text{ref}$ at the load terminals, using a
        proportional-resonant (PR) regulator in the stationary frame
        or a PI regulator with feed-forward in the $dq$ frame.
\end{enumerate}

% --- PLACEHOLDER: control block diagram ---
\begin{figure}[htbp]
  \centering
  \fbox{\parbox{0.9\columnwidth}{\centering\vspace{2cm}%
    [PLACEHOLDER: F-UPQC control block diagram]\vspace{2cm}}}
  \caption{F-UPQC decoupled control architecture.}
  \label{fig:control_overview}
\end{figure}

\subsection{Shunt Inner Current Loop}

The current controller cancels the cross-coupling terms
$\pm\omega_0 L_\text{sh}$ by feed-forward, yielding two independent
first-order plants. A proportional-integral (PI) controller
\begin{equation}
  C_\text{sh}(s) = K_{p,\text{sh}} + \frac{K_{i,\text{sh}}}{s}
  \label{eq:pi_shunt}
\end{equation}
is designed to achieve a closed-loop bandwidth
$\omega_{c,\text{sh}} \approx 2\pi \cdot 1\,\mathrm{kHz}$, well below
the switching frequency.

\subsection{DC Bus Voltage Outer Loop}

The outer voltage loop bandwidth is set one decade below the inner
current loop:
\begin{equation}
  C_\text{dc}(s) = K_{p,\text{dc}} + \frac{K_{i,\text{dc}}}{s},
  \qquad \omega_{c,\text{dc}} \approx 2\pi \cdot 100\,\mathrm{Hz}
  \label{eq:pi_dc}
\end{equation}
This separation of timescales ensures that DC bus dynamics do not
interfere with the inner current regulation.

\subsection{Series Voltage Compensation Loop}

The series controller operates as a voltage feed-forward compensator.
The reference injection voltage is computed as:
\begin{equation}
  \bm{v}_\text{inj}^* = \bm{v}_\text{ref} - \bm{v}_\text{pcc}
                        + Z_\text{line}\,\hat{\bm{i}}_\text{load}
  \label{eq:series_ref}
\end{equation}
where $\hat{\bm{i}}_\text{load}$ is the estimated load current.
A PI regulator with bandwidth $\omega_{c,\text{se}}$ closes the
tracking loop.

% ============================================================
%  V. STEADY-STATE ANALYSIS
% ============================================================

\section{Steady-State Analysis and Component Rating}
\label{sec:steadystate}

\subsection{Power Flow}

Under balanced, steady-state conditions with a purely reactive
ancillary service demand $Q^* \neq 0$ and no active series
compensation ($\bm{v}_\text{inj} = 0$), the active power
balance of the DC bus requires:
\begin{equation}
  P_\text{sh} + P_\text{se} + P_\text{sto} = P_\text{loss}
  \label{eq:power_balance}
\end{equation}
Since reactive power exchange does not transfer net active energy,
$P_\text{sh} \approx P_\text{loss}$, and the storage is idle
($P_\text{sto} \approx 0$) during normal operation.

During a voltage sag of depth $\Delta V$ and duration $\tau$,
the series VSC must inject active power:
\begin{equation}
  P_\text{se,sag} = \frac{\Delta V}{V_\text{nom}}\,P_\text{load}
  \label{eq:sag_power}
\end{equation}
If the shunt VSC cannot fully supply this active power from the grid
(e.g., during a deep or long sag), the energy deficit
\begin{equation}
  E_\text{sto} = P_\text{se,sag}\,\tau
  \label{eq:storage_energy}
\end{equation}
must be supplied by the storage element, which thus determines
the minimum storage capacity for a given sag profile.

\subsection{Converter Rating}

Table~\ref{tab:ratings} summarises the nominal ratings of the main
power components for a representative 100\,kVA installation.

\begin{table}[htbp]
  \centering
  \caption{Nominal Component Ratings (Representative 100\,kVA System)}
  \label{tab:ratings}
  \begin{tabular}{@{}lll@{}}
    \toprule
    \textbf{Component} & \textbf{Parameter} & \textbf{Value} \\
    \midrule
    Shunt VSC           & Apparent power     & $S_\text{sh}$ [kVA] \\
                        & DC link voltage    & $V_\text{dc}$ [V] \\
                        & Coupling inductance & $L_\text{sh}$ [mH] \\
    \midrule
    Series VSC          & Apparent power     & $S_\text{se}$ [kVA] \\
                        & Injection transformer ratio & $n_s$ \\
                        & Series inductance  & $L_\text{se}$ [mH] \\
    \midrule
    DC bus              & Capacitance        & $C_\text{dc}$ [mF] \\
    \midrule
    Storage             & Energy capacity    & $E_\text{sto}$ [kWh] \\
                        & Peak power         & $P_\text{sto}$ [kW] \\
    \bottomrule
  \end{tabular}
\end{table}

% ============================================================
%  VI. REGULATORY CONTEXT
% ============================================================

\section{Regulatory Alignment}
\label{sec:regulatory}

The F-UPQC topology is well-positioned within the evolving European
regulatory framework for active distribution networks.

The \textbf{EU Clean Energy Package}~\cite{ref:CEP} mandates that DSOs
transition from passive network managers to active facilitators of
distributed flexibility. DSO Entity, representing over 900 distribution
operators across 27 EU member states, is developing flexibility
assessment methodologies and market frameworks that require
grid-connected devices to be capable of providing reactive power
support and harmonic mitigation on demand~\cite{ref:EDSO_transition}.

The \textbf{EN~50160} standard defines voltage quality characteristics
at the point of supply, establishing the tolerance bands for voltage
magnitude, harmonic distortion, and unbalance that the series
compensation loop of the F-UPQC is designed to enforce at the load
terminals.

The \textbf{IEC~61000-3-x} series governs harmonic current emission
limits for equipment connected to public networks; the shunt
converter of the F-UPQC can be programmed to ensure compliance of
the load seen from the PCC with these limits, effectively acting
as an active front end (AFE) for the entire facility.

Finally, the \textbf{EU Network Code on Demand Response} (currently
under development alongside Implementing Regulation on data
requirements for demand response~\cite{ref:CEP}) establishes the
conditions under which grid-connected inverters can participate in
distribution-level balancing and flexibility markets. The F-UPQC,
with its bidirectional DC bus and storage interface, is a natural
candidate for participation in such markets, offering a revenue
stream that can partially offset capital expenditure.

% ============================================================
%  VII. CONCLUSION
% ============================================================

\section{Conclusion}
\label{sec:conclusion}

This paper has presented the Flipped UPQC (F-UPQC), a novel topology
for medium-voltage and low-voltage active distribution networks
that simultaneously addresses two emerging requirements: the
protection of sensitive industrial loads from voltage disturbances,
and the provision of ancillary services to the distribution system
operator.

The key innovation of the F-UPQC with respect to the classical UPQC
lies in the topological inversion of the two converters: the shunt VSC
faces the grid, enabling reactive power injection and harmonic
mitigation at the PCC, while the series VSC faces the load, ensuring
clean voltage delivery regardless of grid disturbances.
An optional energy storage element on the shared DC bus extends
compensation capability to sustained sag events and enables
participation in DSO flexibility markets.

A complete mathematical model in the synchronous $dq$ reference frame
has been derived, and a decoupled four-loop control architecture has
been proposed. Steady-state analysis confirms the feasibility of
simultaneous dual-mode operation and provides sizing guidelines for
the main power components.

Future work will focus on:
\begin{itemize}
  \item Hardware-in-the-loop (HIL) validation using a dSPACE
        real-time platform;
  \item Optimisation of storage sizing for participation in
        distribution-level flexibility markets under EN~50160
        compliance constraints;
  \item Extension of the topology to three-phase four-wire systems
        to address voltage unbalance in LV grids.
\end{itemize}

% ============================================================
%  REFERENCES
% ============================================================

\bibliographystyle{IEEEtran}

\begin{thebibliography}{99}

\bibitem{ref:UPQC_original}
H.~Fujita and H.~Akagi,
``The unified power quality conditioner: The integration of series
active filters and shunt active filters,''
\emph{IEEE Trans.\ Power Electron.}, vol.~13, no.~2,
pp.~315--322, Mar.\ 1998.

\bibitem{ref:DVR_original}
N.~H.~Woodley, L.~Morgan, and A.~Sundaram,
``Experience with an inverter-based dynamic voltage restorer,''
\emph{IEEE Trans.\ Power Del.}, vol.~14, no.~3,
pp.~1181--1186, Jul.\ 1999.

\bibitem{ref:STATCOM_ref}
C.~Schauder and H.~Mehta,
``Vector analysis and control of advanced static VAR compensators,''
\emph{IEE Proc.-C}, vol.~140, no.~4, pp.~299--306, Jul.\ 1993.

\bibitem{ref:APF_ref}
B.~Singh, K.~Al-Haddad, and A.~Chandra,
``A review of active filters for power quality improvement,''
\emph{IEEE Trans.\ Ind.\ Electron.}, vol.~46, no.~5,
pp.~960--971, Oct.\ 1999.

\bibitem{ref:EDSO_transition}
DSO Entity,
``DSOs in the Energy Transition,''
[Online]. Available: \url{https://eudsoentity.eu},
Accessed: 2025.

\bibitem{ref:EC_grids}
European Commission,
``Grids, the Missing Link -- An EU Action Plan for Grids,''
COM(2023) 757 final, Brussels, Nov.\ 2023.

\bibitem{ref:ancillary_EU}
M.~Marinelli \emph{et al.},
``Ancillary services markets in Europe: Evolution and regulatory
trade-offs,''
\emph{Renew.\ Sustain.\ Energy Rev.}, vol.~151, p.~111566, 2021.

\bibitem{ref:CEP}
European Parliament and Council,
``Clean Energy for All Europeans Package,''
Directives (EU) 2018/2001, 2019/944, Brussels, 2018--2019.

\bibitem{ref:IEEE_std_519}
IEEE Std 519-2022,
\emph{IEEE Standard for Harmonic Control in Electric Power Systems},
IEEE, New York, 2022.

\bibitem{ref:PQ_cost}
Fortune Business Insights,
``Power Quality Equipment Market,''
Market Research Report, 2025.
[Online]. Available: \url{https://www.fortunebusinessinsights.com/power-quality-equipment-market-111978}

\end{thebibliography}

\end{document}
